\section{Conclusion}
\label{sec:conclusion}

We presented the first empirical study specifically targeting thankfulness as an LLM behavior, testing gratitude along three dimensions: expression, reception, and behavioral grounding. Our experiments with \gptfour reveal a consistent pattern: the model has learned sophisticated contextual rules for \emph{when} to express gratitude (95\% accuracy) and \emph{mirrors} user gratitude by producing longer responses, but self-reported gratitude does not predict cooperative behavior. LLM thankfulness is best understood as a linguistically competent performance shaped by RLHF training rather than evidence of genuine preference states.

The key takeaway for practitioners is twofold. First, LLM gratitude is functionally useful---context-appropriate expressions of thanks make interactions feel natural---even if it is not genuine. Second, self-report measures of LLM traits should not be trusted as indicators of behavioral tendencies, extending prior findings on personality to the specific domain of gratitude.

Future work should extend this investigation to multiple model families and scales, employ mechanistic interpretability to probe whether gratitude-related representations exist in internal activations, and design multi-turn interactive scenarios that test behavioral cooperation in richer settings. The question that motivated this work---how should we design the appearance of AI preferences to interact with people best?---remains the right question to ask, and our data provides an empirical foundation for answering it.